\section{Discussion}
\label{sec:discussion}

\subsection{Why Local-First Is the Right Default}
The case study completed a static, semantic, structural, and dynamic
audit with no code leaving the machine. For the teams described in
\cref{sec:intro}---proprietary, regulated, or air-gapped---this is the
difference between a tool they can run today and one they cannot run at
all. The cost is borne entirely in model capability, which we consider a
better trade than the alternative, because the pattern and complexity
passes are unaffected by model size and the model's role is advisory and
always shown next to corroborating evidence.

\subsection{The Recall Cost of a Small Model}
A quantised \mbox{Llama~3.2} model is far smaller than the frontier models
used in most published LLM-for-security
work~\cite{pearce2023,khoury2023}. We expect it to miss subtle,
cross-function, or dataflow-dependent vulnerabilities that a larger model
or a dedicated learned detector~\cite{fu2022linevul} would catch, and to
occasionally produce plausible but incorrect findings. \cs\ mitigates
this in three ways: it constrains the model to a strict output schema and
discards malformed items; it fuses model output with deterministic rules
so that the unambiguous classes do not depend on the model; and it never
auto-applies a fix. The interface is explicit that model findings are
suggestions for a human to confirm.

\subsection{Buildability and Dynamic Coverage}
The dynamic pass requires a project that builds. When a repository has no
Dockerfile, \cs\ synthesises a minimal one, which works for
conventionally structured projects but will fail for those with unusual
build systems or external service dependencies. The current behavioural
event stream is coarse and, in the prototype, partly simulated for
demonstration; a production version would attach a real syscall or eBPF
monitor. Dynamic analysis in \cs\ is therefore a fast liveness-and-shape
check, not a substitute for fuzzing~\cite{manes2021}.

\subsection{Threats to Validity}
\textbf{Construct.} The \phindex\ weights in \cref{eq:h} are chosen by
design intent, not fitted to outcome data; the score is best read as an
ordinal triage aid, and its calibration against expert judgement or
incident data is open work. \textbf{Internal.} Finding fusion keys on
$(\textit{file}, \textit{line}, \textit{category})$; line drift between
the pattern and model passes can under- or over-merge, though severity is
always taken as the maximum. \textbf{External.} The evaluation is a
single, synthetic subject chosen to exercise a specific weakness class.
It demonstrates what an audit produces and how it is presented; it does
not establish precision, recall, or false-positive rate on real-world
code. A study on a labelled corpus (with the dataset-quality caveats of
\cite{croft2023,steenhoek2023}) and a comparison against established
tools is required before efficacy claims can be made. \textbf{Conclusion.}
Runtime figures are from one commodity machine and one workload and will
vary with hardware and model quantisation.

\subsection{Future Work}
Four directions follow directly. \emph{(i)~Empirical evaluation} on a
labelled multi-language vulnerability corpus, reporting precision/recall
against ground truth and inter-tool agreement. \emph{(ii)~Deeper static
analysis}: replace regex rules with an AST/code-property-graph
layer~\cite{yamaguchi2014} to add inter-procedural taint tracking.
\emph{(iii)~Model scaling and routing}: support larger local models where
hardware permits, and route only ambiguous files to the model to control
latency. \emph{(iv)~Workflow integration}: a headless mode and a
pre-commit / CI gate driven by a \phindex\ threshold, plus a controlled
user study of triage speed against the usability criteria of
\cite{johnson2013,nachtigall2022}.
